\documentclass[10pt]{article}
\usepackage[T1]{fontenc}
\usepackage[utf8]{inputenc}
\usepackage[french]{babel}
\usepackage{amsmath,amssymb,mathrsfs,mathpazo}
\usepackage{geometry}
\geometry{a4paper,margin=2cm}
\usepackage{graphicx,array,tabularx,tabularray,colortbl,multirow,multicol,booktabs,enumitem,pifont,ulem,fancybox,diagbox,cancel}
\usepackage{pstricks,pstricks-add,pst-plot,pst-tree,pst-node,pst-text,pst-func,pst-math,pst-bspline,pst-3dplot,pst-eucl,pst-solides3d}
\let\pscalculate\undefined
\usepackage{pst-bezier}
\usepackage{tikz,pgfplots,tkz-tab}
\usetikzlibrary{arrows,arrows.meta,patterns,calc,positioning,trees,decorations.markings,shapes.geometric,shapes.misc,shapes.arrows,intersections,decorations.pathmorphing,decorations.pathreplacing,backgrounds,shadows,babel}
\pgfplotsset{compat=1.18}
\usepgfplotslibrary{fillbetween}
\usepackage[np]{numprint}
\usepackage{listings,makecell,pas-tableur,twoopt}
\let\barre\undefined
\usepackage[pstricks]{tablvar}
\definecolor{darkgreen}{RGB}{0,100,0}
\providecommand{\red}{\color{red}}
\providecommand{\blue}{\color{blue}}
\providecommand{\green}{\color{green}}
\providecommand{\black}{\color{black}}
\providecommand{\white}{\color{white}}
\providecommand{\gray}{\color{gray}}
\providecommand{\cyan}{\color{cyan}}
\providecommand{\magenta}{\color{magenta}}
\providecommand{\yellow}{\color{yellow}}
\providecommand{\violet}{\color{violet}}
\providecommand{\orange}{\color{orange}}
\providecommand{\euro}{\texteuro}
\providecommand{\EUR}{\texteuro}
\providecommand{\ieme}{\textsuperscript{e}}
\providecommand{\degres}{\ensuremath{^\circ}}
\providecommand{\degree}{\ensuremath{^\circ}}
\providecommand{\vv}[1]{\overrightarrow{#1}}
\providecommand{\Coord}[1]{#1}
\providecommand{\decosix}{}
\providecommand{\starredbullet}{\ensuremath{\bullet}}
\pagestyle{empty}
\setlength{\parindent}{0pt}
\usepackage{adjustbox}\DeclareUnicodeCharacter{2236}{:}\DeclareUnicodeCharacter{2190}{\ensuremath{\leftarrow}}
\providecommand{\R}{\mathbb{R}}
\providecommand{\N}{\mathbb{N}}
\providecommand{\Z}{\mathbb{Z}}
\providecommand{\Q}{\mathbb{Q}}
\providecommand{\D}{\mathbb{D}}
\providecommand{\C}{\mathbb{C}}
\providecommand{\e}{{\mathrm{e}}}
\providecommand{\vect}[1]{\overrightarrow{#1}}
\providecommand{\vectt}[1]{\overrightarrow{\mathrm{#1}}}
\providecommand{\vvt}[1]{\overrightarrow{\mathrm{#1}}}
\providecommand{\np}[2][]{{#2~#1}}
\providecommand{\vv}[1]{{\overrightarrow{#1}}}
\renewcommand{\d}{{\mathrm{d}}}
\providecommand{\degres}{{^\circ}}
\providecommand{\degre}{{^\circ}}
\providecommand{\blue}{\color{blue}}
\providecommand{\red}{\color{red}}
\providecommand{\enskip}{\quad}
\providecommand{\llbracket}{\lbrack\!\lbrack}
\providecommand{\rrbracket}{\rbrack\!\rbrack}
\providecommand{\euro}{{\text{€}}}
\providecommand{\no}{{\text{n°}}}
\providecommand{\interval}[2]{{[#1;#2]}}
\providecommand{\intervFO}[2]{{[#1;#2[}}
\providecommand{\dotfill}{\ldots}

\providecommand{\ds}{\displaystyle}
\providecommand{\textding}[1]{#1}
\providecommand{\Oij}{\(\mathrm{O};\vec{\imath},\vec{\jmath}\)}
\providecommand{\Oijk}{\(\mathrm{O};\vec{\imath},\vec{\jmath},\vec{k}\)}
\providecommand{\Ouv}{\(\mathrm{O};\vec{u},\vec{v}\)}
\providecommand{\intervFF}[2]{\ensuremath{[#1;#2]}}
\renewcommand{\euro}{\ifmmode\text{\texteuro}\else\texteuro\fi}
\DeclareUnicodeCharacter{00B0}{\ifmmode^{\circ}\else\textdegree\fi}
\DeclareUnicodeCharacter{20AC}{\ifmmode\text{\texteuro}\else\texteuro\fi}
% Énoncés : Éducation nationale. Transcription et corrigés : APMEP, Denis Vergès et contributeurs. Découpage et classement : CoopMaths, Sébastien Lozano, Éric Elter et contributeurs. Adaptation de la présentation : Maths Explorer.
% CC BY-SA : https://coopmaths.fr/www/about/
% https://coopmaths.fr/alea/static/bac/2024/tex/bac_2024_06_sujet2_etrangers_4.tex
\begin{document}
\noindent\textit{APMEP / CoopMaths. CC BY-SA. Présentation adaptée par Maths Explorer.}\par\medskip
\medskip

\textbf{Partie A}

\medskip

On considère la fonction $f$ définie sur l'intervalle $[0~;~+\infty[$ par 
\[f(x) = \sqrt{x + 1}.\]

On admet que cette fonction est dérivable sur ce même intervalle.

\medskip

\begin{enumerate}
\item Démontrer que la fonction $f$ est croissante sur l'intervalle $[0~;~+\infty[$.
\item Démontrer que pour tout nombre réel $x$ appartenant à l'intervalle $[0~;~+\infty[$ :

\[f(x) - x = \dfrac{-x^2 + x + 1}{\sqrt{x + 1} + x}.\]

\item En déduire que sur l'intervalle $[0~;~+\infty[$ l'équation $f(x) = x$ admet pour unique solution :

\[\ell = \dfrac{1 +\sqrt 5}{2}.\]

\end{enumerate}

\medskip

\textbf{Partie B}

\medskip

On considère la suite $\left(u_n\right)$ définie par $u_0 = 5$ et pour tout entier naturel $n$, par 

$u_{n+1} = f\left(u_n\right)$ où $f$ est la fonction étudiée dans la \textbf{partie A}.

On admet que la suite de terme général $u_n$ est bien définie pour tout entier naturel $n$.

\medskip

\begin{enumerate}
\item Démontrer par récurrence que pour tout entier naturel $n$, on a 

\[1 \leqslant u_{n+1} \leqslant u_n.\]

\item En déduire que la suite $\left(u_n\right)$ converge.
\item Démontrer que la suite $\left(u_n\right)$ converge vers $\ell = \dfrac{1 +\sqrt 5}{2}$.
\item On considère le script Python ci-dessous :
\begin{center}
\begin{tabular}{|>{\ttfamily}l >{\ttfamily}l|}\hline
1 &\textbf{from} math \textbf{import} *\\
2 &\textbf{def} seuil(n):\\
3 &\qquad u=5\\
4 &\qquad i=$0$\\
5 &\qquad $\ell$ =(1 + \textbf{sqrt}(5))/2\\
6 &\qquad \textbf{while abs}(u-$\ell$)>=10**(-n):\\
7 &\qquad \qquad u=\textbf{sqrt}(u+1)\\
8 &\qquad \qquad i=i+1\\
9 &\qquad \textbf{return}(i)\\ \hline
\end{tabular}
\end{center}

\emph{On rappelle que la commande} \textbf{abs(x)} \emph{renvoie la valeur absolue de $x$}.

	\begin{enumerate}
		\item Donner la valeur renvoyée par {\ttfamily seuil} (2).
		\item La valeur renvoyée par {\ttfamily seuil} (4) est 9.
		
Interpréter cette valeur dans le contexte de l'exercice.
	\end{enumerate}
\end{enumerate}

\end{document}
